% ===== c07：附录 C 四维的 $\eta$ 不变量（The $\eta$-Invariant In Four Dimensions）与参考文献 =====
% 源文件 QuantumWorkTry.tex 第 2919--3242 行（对应 fermi/chunks/c07.tex）
% 说明：本文件仅含中文译注正文片段，不含导言区与 \end{document}

\section{四维的 $\eta$ 不变量（The \ensuremath{\eta}-Invariant In Four Dimensions）}\label{feta}

\begin{keytakeaways}
\begin{itemize}[leftmargin=1.6em,itemsep=0.2em,topsep=0.2em]
  \item \textbf{物理图像}：四维不可定向流形上，pin 结构的马约拉纳费米子路径积分由一个谱不对称相位给出；该相位只能取十六次单位根，这正是三维拓扑超导体十六重分类在四维体相一侧的场论根源。
  \item \textbf{数学工具}：APS 指标定理的奇维形式（指标等于负的半个谱不对称不变量）、定向双重覆盖、相差定向线丛的互补 pin 结构，以及实向量丛的 Stieffel--Whitney 类与第一庞特里亚金类。
  \item \textbf{核心结论}：两个互补 pin 结构的谱不对称不变量之和模 4 为零，其八倍之差模 4 亦为零，故该不变量必为四分之一的整数倍；实射影四维空间上它等于负四分之一（符号依赖 pin 结构之选取），从而相位恰为十六次单位根。
\end{itemize}
\end{keytakeaways}

在 \ref{bands} 节与 \ref{zeb} 节中，我们需要一些关于 $D=4$ 中 $\eta$ 不变量的事实。\ref{maj} 节中则援引了 $D=2$ 中的类似事实。本附录的目标是对这些事实作一简要说明。（我们给出的论证与最初的论证 \cite{G} 相近。）

我们的第一个目标是证明：在 $D=4$ 中，$\eta$ 总是 $1/4$ 的整数倍，从而 $\exp(-\i\pi \eta/2)$ 总是 1 的 16 次单位根。非平凡的情形是 $X$ 不可定向；否则 $\eta$ 是 2 的整数倍，正如 \ref{bpf} 节末尾所指出。设 $X$ 具有 $\pin^+$ 结构 $\P$ 以及互补的 $\pin^+$ 结构 $\P'=\P\otimes \veps$，其中 $\veps $ 是一个实线丛，即 $X$ 的定向丛。对于取值于这两种 $\pin^+$ 结构中任何一种的马约拉纳费米子，我们都可以定义 $\eta$ 不变量 $\eta_\P$、$\eta_{\P'}$。
% [Tutor Note: 关键约定在于 $\veps$：它是 $X$ 的定向丛，故 $w_1(\veps)=w_1(TX)$，其非零当且仅当 $X$ 不可定向；$\P'=\P\otimes\veps$ 是把 $\P$ 的转移函数再乘以定向丛的符号而得的另一个 $\pin^+$ 结构。因实线丛满足 $\veps\otimes\veps\cong\R$，故 $\P'\otimes\veps=\P$，二者互为互补结构；不可定向四维流形上的两个 $\pin^+$ 结构恰是这样一对，$\eta_\P$ 与 $\eta_{\P'}$ 是仅有的两个可定义的 $\eta$ 不变量。]

第一个有用的关系是
\be\label{simpler}\eta_\P+\eta_{\P'}=0 ~~~{\mathrm {mod}}~4. \ee
% [Tutor Note: 式（\ref{simpler}）的证明要点是"直和即提升"：把分别耦合于 $\P$ 与 $\P'$ 的两份马约拉纳费米子直和起来，所得正是 $\h X$ 上耦合于某个自旋结构的单份马约拉纳费米子。理由是 $\veps$ 提升到定向双重覆盖 $\h X$ 上成为平凡丛，于是 $\P$ 与 $\P'$ 的提升给出同一个自旋结构，而两者的截面分别对应在 Deck 变换下为偶、为奇的旋量场，合起来即 $\h X$ 的全部旋量场。因此 $\eta_\P+\eta_{\P'}=\h\eta$ 是精确等式而非同余，取模 4 只是后续推论；务必不要把这一步误当作只差一个 4 的整数倍。]

要点在于：$X$ 上耦合于 $\P$ 的马约拉纳费米子，与 $X$ 上耦合于 $\P'$ 的马约拉纳费米子合在一起，构成 $X$ 的定向双重覆盖 $\h X$ 上一个耦合于自旋结构的马约拉纳费米子。因此 $\eta_\P+\eta_{\P'}=\h\eta$，其中 $\h\eta$ 是 $\h X$ 上马约拉纳费米子的 $\eta$ 不变量。由于 $\h X$ 可定向，并且具有一个反定向的对称性（其商为 $X$），故 $\h\eta$ 是 4 的整数倍，正如 \ref{bpf} 节末尾所指出。
% [Tutor Note: "$\h\eta$ 是 4 的整数倍"的原始论证在 \ref{bpf} 节末：对可定向流形，非零本征值成对 $\lambda\leftrightarrow-\lambda$ 出现而不贡献 $\eta$，故 $\eta$ 只由零模给出，$\eta=n_++n_-$；赝实性使每个本征值（含零模）具有偶重数，故 $n_+$、$n_-$ 均为偶数。$\h X$ 上还存在反定向对合（商为 $X$），它交换正、负手征，给出 $n_+=n_-$，于是 $\eta=2n_+$ 被 4 整除。这正是"$\h X$ 的 $\eta$ 为什么是 4 的倍数"的全部内容。]

我们还需要另一个关系：
\be\label{impler}8\eta_{\P'}=8\eta_\P~~{\mathrm{mod}}~4. \ee
% [Tutor Note: 由这两个关系推出 $\eta_\P$、$\eta_{\P'}$ 均为 $1/4$ 的整数倍，原文省略了一步初等数论。记 $a=\eta_\P$、$b=\eta_{\P'}$。第一式给出 $a+b=4k$（$k\in\Z$）；第二式即 $8(a-b)$ 被 4 整除，故 $a-b\in\frac12\Z$。于是 $a=\frac{a+b}{2}+\frac{a-b}{2}=2k+\frac{a-b}{2}$，其中 $2k\in\Z$、$\frac{a-b}{2}\in\frac14\Z$，故 $a\in\frac14\Z$；同法得 $b\in\frac14\Z$。仅凭 $a+b\in4\Z$ 只能得到 $a,b\in\frac12\Z$，必须配合第二式才能收紧到 $1/4$ 的整数倍。]
这两个关系合起来表明，$\eta_\P$ 与 $\eta_{\P'}$ 都是 $1/4$ 的整数倍。

为解释式（\ref{impler}），把问题稍作推广是有益的。设 $V\to X$ 为任意实向量丛。我们考虑作用在取值于 $V$ 的费米子上的狄拉克算子 $\D_V$，即作用在 $\P\otimes V$ 的截面上的算子。对任意 $V$，我们都可以定义 $\eta_{\P\otimes V}$，即作用在这类场上的狄拉克算子的 $\eta$ 不变量。这里，对任意 $V$，$\eta_{\P\otimes V}$ 是模 4 的拓扑不变量（与 $X$ 的度规以及 $V$ 上的联络无关）。这可由 APS 指标定理推出，该定理在奇数维情形下取简单形式（\ref{molg}）。为在不必假设 $X$ 是边界的前提下应用该定理，我们考虑 5 维流形 $Z=X\times [0,1]$，使其一端具有度规与规范场 $(g,A)$，另一端具有另一组 $(g',A')$。该定理指出，两端 $-\eta/2$ 之差等于 $\I$，即 $Z$ 上带 APS 边界条件的狄拉克算子的指标。如 \ref{fvod} 节所解释，该指标总是偶数，因此两个 $\eta$ 值之差是 4 的整数倍。更直接地说，实际发生的情形是：当我们在偶数维流形上改变 $(g,A)$ 时，$\eta$ 保持不变，唯一的例外是当某个本征值穿过 0 时它跳变 2。在 4 维中，由于某种形式的 Kramers 二重化，本征值具有偶重数，因此这些跳变是 4 的倍数，而 $\eta$ 是模 4 的拓扑不变量。
% [Tutor Note: 这一段省略了 APS 指标定理在奇维的完整陈述与用法。把两端的"跃变"实现为柱体 $Z=X\times[0,1]$ 上的指标，是 APS 的标准技巧：$Z$ 上带 APS 边界条件（只保留边界上负本征值对应的模式）的狄拉克指标等于两端 $-\eta/2$ 之差；若两端配同样的几何，指标为零即得 $\eta$ 的形变不变性。该指标恒为偶数，来源是五维的赝实性------如 \ref{fvod} 节所示，$\Pin^+(5)$ 的旋量表示上存在反线性算子 $T$（$T^2=-1$），它保证 $\P$ 与 $\P'$ 上的零模数各自为偶。后一论证（本征值穿越零时的跳变）与此等价：单个本征值由负穿到正使 $\eta$ 改变 2，而四维中本征值因 Kramers 型二重化具有偶重数，故实际跳变是 4 的倍数。]
% [Tutor Note: 原文本段有两处排印疏漏："is helpful to generalize"之前脱漏形式主语 it；末句"jumps are multiplies of 4"中的 multiplies 应为 multiples。译文按原意处理，公式与文字均未改动。]

我们可以把 $8\eta_{\P'}$ 解释为 $\eta_{\P\otimes V}$，其中 $V=\veps^{\oplus 8}$ 是实线丛 $\veps$ 的 8 个副本的直和。类似地，$8\eta_\P=\eta_{\P\otimes \R^8}$，其中 $\R^8$ 是秩为 8 的平凡实向量丛。因此，为建立式（\ref{impler}），只需证明 $\veps^{\oplus 8}$ 是平凡的，即同构于 $\R^8$。

为证明 4 维流形 $X$ 上的实向量丛 $V$ 平凡，只需证明 Stieffel-Whitney 类 $w_i(V)$（$1\leq i\leq 4$）为零，并证明某个类 $\lambda\in H^4(X,\Z)$ 为零；后者的定义颇为微妙，但它满足 $2\lambda=p_1(V)$，其中 $p_1$ 是第一庞特里亚金类。
% [Tutor Note: 用 $w_i$（$i=1,\dots,4$）与 $\lambda$ 判别四维实丛的平凡性，属于四维实丛的障碍理论：$w_1=w_2=0$ 给出自旋（或 pin）提升条件，$w_3$、$w_4$ 与 $\lambda$ 则控制剩余障碍。$\lambda$ 之所以"微妙"，是因为一般实丛的 $p_1$ 未必是偶类；对满足 $w_1=w_2=0$ 的实丛，因 $p_1\equiv w_2^2\pmod 2$，可证明 $p_1=2\lambda$ 且 $\lambda$ 为整上同调类。原文只对秩 8 的 $\veps^{\oplus 8}$ 逐项验证这些障碍为零，未展开一般理论。]

关于第一点，$V$ 的总 Stieffel-Whitney 类为 $w(V)=w(\veps^{\oplus 8})=w(\veps)^8=(1+w_1(\veps))^8=1$，其中最后一步用到相关的二项式系数模 2 为零这一事实。于是对 $i\geq 1$ 有 $w_i(V)=0$。
% [Tutor Note: 这一步值得展开：在 $\Z_2$ 系数下 Whitney 乘积公式给出 $w(\veps^{\oplus 8})=w(\veps)^8=(1+w_1(\veps))^8$；由特征 2 的"新生之梦"$(a+b)^8=a^8+b^8$ 得 $1+w_1^8$，而 $w_1^8\in H^8(X,\Z_2)=0$（$X$ 为四维）。逐项看，二项式系数 $\binom{8}{k}$（$k=1,\dots,7$）均为偶数，故中间各项同样消失。顺带指出：原文"Stieffel-Whitney"是 Stiefel--Whitney（施蒂费尔--惠特尼）的拼写笔误，译文保留原拼写。]

为证明 $\lambda(V)$ 为零，我们先把该丛写成 $V=V_0\oplus V_0$，其中 $V_0=\veps^{\oplus 4}$。于是 $\lambda(V) =2\lambda(V_0)=p_1(V_0)$。这样问题就归结为证明 $p_1(V_0)=0$。为计算 $p_1(V_0)$，我们注意到 $V_0$ 的复化是 $\veps_\C=\veps\otimes_\R\C$ 的 4 个副本的直和。一般地，若 $V_0\to X$ 是实向量丛，其复化为若干复线丛 $\L_i\to X$ 的直和 $\oplus_i \L_i$，则 $p_1(V_0)=\sum_i c_1(\L_i)^2$。在目前情形下，该公式给出 $p_1(V_0)=4c_1(\veps_\C)^2$，而它为零，因为 $c_1(\veps_\C)$ 是 2-挠的。最后这一论断是普遍成立的：对拓扑空间上的任意实线丛 $\veps$，其复化 $\veps_\C=\veps\otimes_\R\C$ 均满足该论断。
% [Tutor Note: "$c_1(\veps_\C)$ 是 2-挠"只需一句证明：实线丛的复化带有实结构，复共轭给出复线丛同构 $\overline{\veps_\C}\cong\veps_\C$；而 $c_1(\bar E)=-c_1(E)$，故 $c_1(\veps_\C)=-c_1(\veps_\C)$，即 $2c_1(\veps_\C)=0$，于是 $4c_1(\veps_\C)^2=0$。对本题的四维 $X$，这一点也可直接从 $H^2(X,\Z)$ 的 2-挠性看出（如 $\RP^4$ 的 $H^2=\Z_2$）。]

至此便完成了 $\exp(-\i\pi \eta/2)$ 是 1 的 16 次单位根这一证明。
% [Tutor Note: 这条推导链可概括为：$\eta_\P\in\frac14\Z$（由前两个关系）$\Rightarrow$ $\eta_\P/2\in\frac18\Z$ $\Rightarrow$ $\exp(-\i\pi\eta_\P/2)=\exp(-2\pi\i\,m/16)$（$m\in\Z$），即十六次单位根。指数中出现 $\eta/2$ 而非 $\eta$，是马约拉纳费米子配分函数取 Pfaffian、需从每一对里只取一个本征值之故（参见 \ref{bpf} 节的 $\prod_k'$ 约定）；这也解释了为什么这里要求 $\eta$ 模 4 而不是模 2。]
% [Tutor Note: 接下来这一段是本节唯一的显式计算，思路是"以边界算体"：直接在 $\RP^4$ 上算 $\eta$ 并不方便，于是把 $\RP^4$ 作为五维带边流形 $X$ 的边界，由五维 APS 公式 $\eta/2=-\I$ 把 $\eta$ 化为指标；该指标又可在覆盖空间 $T^5$（挖去 32 个小球）上用常旋量直接数出。]

我们还希望证明：对 $X=\RP^4$，有 $\exp(-\pi i\eta/2)=\exp(\pm 2\i\pi /16)$，其中符号取决于 $\Pin^+$ 结构的选取。我们从五维环面 $T^5$ 出发，其参数为周期为 1 的变量 $x_1,\dots,x_5$。我们给 $T^5$ 一个自旋结构，使得旋量在所有方向上都是周期的。$T^5$ 的旋量丛 $\S$ 秩为 4，因为 $\SO(5)$ 的旋量表示是 4 维的。

$T^5$ 上的狄拉克算子 $\D=\i \slashed{\partial}$ 是自伴的，故其指标显然为零。现在我们用 $T^5/\Z_2$ 代替 $T^5$，其中 $\Z_2$ 通过 $\sP:x_i\to -x_i$、$i=1,\dots,5$ 作用。这里 $T^5/\Z_2$ 并不是流形，我们稍后将对之作修正。我们定义 $\H_+$ 与 $\H_-$ 为 $T^5$ 上分别满足 $\psi(-x)=\psi(x) $ 与 $\psi(-x)=-\psi(-x)$ 的旋量场空间。
% [Tutor Note: 原文"$\psi(-x)=-\psi(-x)$"疑为笔误：按上下文（$\H_-$ 为奇场空间）应为 $\psi(-x)=-\psi(x)$。译文保留原式，以与源码逐字对应。]
由于 $\D$ 与 $\sP$ 反对易，它把 $\H_+$ 映到 $\H_-$，反之亦然。我们把 $\D$ 在 $T^5/\Z_2$ 上的\emph{指标} $\I$ 定义为：$\D$ 作用在 $\H_+$ 上的零模个数，减去 $\D$ 作用在 $\H_-$ 上的零模个数。（这两个算子互为伴随，因为作用在 $\H_+\oplus \H_-$ 上的完整狄拉克算子是自伴的。）

$\I$ 的计算是直接的。$\D$ 作用在 $\H_+$ 上的零模构成 4 维的常旋量空间，而 $\D$ 作用在 $\H_-$ 上没有零模。故指标为 $\I=4$。
% [Tutor Note: $\I=4$ 的计数理由：周期自旋结构下 $\D=\i\slashed{\partial}$ 的本征值为 $\pm 2\pi n$（$n$ 取整格），故零模只能是常旋量；常旋量在 $\sP:x\mapsto -x$ 下为偶，属于 $\H_+$，其空间维数等于 $\Spin(5)\cong\Sp(4)$ 旋量表示的维数 4，而 $\H_-$ 中没有常旋量。注意这里"4"来自旋量表示的维数，不是来自流形的维数。]

现在，$T^5/\Z_2$ 不是流形，因为 $\sP$ 有 32 个不动点：即坐标满足 $x_i\in \Z/2$、$i=1,\dots,5$ 的那些点。为得到流形，我们从 $T^5$ 中挖去每个不动点周围的一个小球，得到带边界的流形 $T'$。$T'$ 对 $\sP$ 的商则是一个带边界的流形 $X$，其边界由 32 个 $\RP^4$ 的副本组成。$X$ 获得两个 $\pin^+$ 结构 $\P$、$\P'$，其截面分别是 $T'$ 上在 $\sP$ 下为偶或为奇的旋量场。狄拉克算子 $\D$ 把 $\P$ 的截面映为 $\P'$ 的截面，反之亦然。作为这样的算子，它具有一个指标 $\I$。一个相当初等的论证\footnote{在 $\sP$ 作用于 $T^5$ 的某个不动点附近，我们可以用欧氏度规为 $\d \ell^2=\d \vec x^2$ 的 $\R^5/\Z_2$ 代替 $T^5/\Z_2$。作一次 Weyl 重标度，可将它变为 $\d\h\ell^2=\d \vec x^2\cdot \left(1+{1}/{|\vec x|^2}\right)$。在 $\vec x=0$ 附近，这具有我们在引入 APS 边界条件时所描述的管状行为（见 \ref{aps} 节的图 \ref{longer}）。在此情形下，该管是 $\RP^4\times \R_+$ 的一个副本。由于狄拉克方程的共形不变性，$\R^5$ 上在 $\vec x=0$ 处正则的 $\Z_2$ 不变解，等价于在管状几何中平方可积的解。因此，我们既可以按原始 $T^5$ 几何中的正则解来计数指标 $\I$，也可以先作 Weyl 变换把 32 个不动点替换为管，再按平方可积解来计数。然而，APS 边界条件的定义保证：后一做法等价于在 $X $ 上带 APS 边界条件计算的指标。}利用狄拉克算子的共形不变性表明：$X$ 上带 APS 边界条件的狄拉克指标，与上一段在 $T^5/\Z_2$ 上求得的 $\I=4$ 相同。
% [Tutor Note: 脚注的论证可概括为"共形泡"（bubbling）：在不动点附近把 $T^5/\Z_2$ 局部换为 $\R^5/\Z_2$，再作 Weyl 变换 $\d\ell^2\mapsto(1+1/|\vec x|^2)\d\vec x^2$，把锥点拉成长管 $\RP^4\times\R_+$；因无质量狄拉克方程在 Weyl 变换下共形不变（解只差一个共形因子），原点处正则的 $\Z_2$ 不变解与管中平方可积解一一对应，而管末端的平方可积条件恰是 APS 边界条件。这正是"在 $T^5/\Z_2$ 上数正则解"与"在 $X$ 上算 APS 指标"相等的理由。]

在奇数维中，APS 指标定理取式（\ref{molg}）的形式，它只是说 $\eta/2=-\I$，其中 $\I$ 是流形 $T'$ 上的狄拉克指标，而 $\eta$ 是其边界的 $\eta$ 不变量。在目前情形下，边界由 32 个 $\RP^4$ 组成（它们都由明显的对称性相互置换，因而具有相同的 $\pin^+$ 结构）。于是公式变为 $16\eta_{\RP^4}=-\I=-4$，即 $\eta_{\RP^4}=-1/4$。
% [Tutor Note: 此处记账易错，务必区分"边界副本数"与"单个 $\RP^4$ 的 $\eta$"：边界有 32 个 $\RP^4$ 副本，故边界整体的 $\eta=32\eta_{\RP^4}$；代入奇维 APS 公式 $\eta/2=-\I$ 得 $16\eta_{\RP^4}=-\I=-4$，从而 $\eta_{\RP^4}=-1/4$。若误把边界整体的 $\eta$ 当作单个 $\eta_{\RP^4}$，会得到 $-1/2$ 的错误结果。符号 $\pm$ 对应两个 $\pin^+$ 结构的选取，即 $\exp(\pm 2\i\pi/16)$ 的来源。]

我们留请读者用类似方法证明：在 2 维中，$\exp(-\i\pi \eta/2)$ 是 1 的 8 次单位根，且 $\eta_{\RP^2}=\pm 1/2$。
% [Tutor Note: 二维情形的推导与本节完全平行：$\RP^2$ 的定向双重覆盖是 $S^2$，其 $\eta$ 为 4 的倍数，由此得 $\eta_{\RP^2}=\pm 1/2$，于是 $\exp(-\i\pi\eta/2)$ 只取八次单位根。维数由 2 升到 4，单位根次数相应由 8 变为 16，正对应拓扑超导体的八重与十六重分类（参见 \cite{KitTwo} 与 \cite{MFCV}）。]

本研究部分受 NSF 基金 PHY-1314311 资助。作者感谢 M. F. Atiyah、G. Burdman、D. Freed、A. Kitaev、S. Ryu 与 N. Seiberg 的评论与讨论。

\bibliographystyle{unsrt}
\begin{thebibliography}{99}
\bibitem{WenEtAl}X. Chen, Z.-C. Gu, Z.-X. Liu, and X.-G. Wen,
``Symmetry Protected Topological Orders And The Group
Cohomology Of Their Symmetry Group,'' Phys. Rev. {\bf B87} (2013) 155114.

\bibitem{CH}
C. G. Callan, Jr., and J. A. Harvey, ``Anomalies And Fermion Zero-Modes On Strings And Domain Walls,''
 Nucl.Phys. {\bf B250} (1985) 427.

\bibitem{FreedTwo}
D. S. Freed, ``Anomalies And Invertible Field Theories,'' arXiv:1404.7224.

\bibitem{Adler}
S. Adler, ``Axial Vector Vertex In Spinor Electrodynamics,'' Phys. Rev. {\bf 177} (1969) 2426-2438.



\bibitem{BJ}
J. S. Bell and R. Jackiw, ``A PCAC Puzzle: $\pi^0\to\gamma\gamma$ In The Sigma Model,''
Nuovo Cimento {\bf A60} (1969) 47-61.

\bibitem{KH}
M. Z. Hasan and C. L. Kane, ``Topological Insulators,'' Rev. Mod. Phys. {\bf 82} (2010) 3045.

\bibitem{ZQ}
Xiao-Liang Qi and Shou-Cheng Zhang, ``Topological Insulators And Superconductors,'' Rev. Mod. Phys. {\bf 83}
(2011)  1057-1110, arXiv:1008.2026.

\bibitem{MZ}
M. Z. Hasan and J. E. Moore, ``Three-Dimensional Topological Insulators,'' arXiv:1011.5462,
Annu. Rev. Cond. Matt. Phys. {\bf 2} (20110 55.


\bibitem{RML}
S. Ryu, J. E. Moore, and A. W. W. Ludwig, ``Electromagnetic And Gravitational Responses And Anomalies In Topological
Insulators And Superconductors,'' arXiv:1010.0936.

\bibitem{Kitaev}
A. Kitaev, ``Periodic Table For Topological Insulators And Superconductors,'' arXiv:0901.2686.

\bibitem{RSFL}
S. Ryu, A. P. Schnyder, A. Furusaki, and A. W. W. Ludwig, ``Topological Insulators And Superconductors:
Tenfold Way And Dimensional Hierarchy,'' New J. Phys. 12, 065010 (2010).



\bibitem{ASfirst}
M. F. Atiyah and I. M. Singer, ``The Index Of Elliptic Operators On Compact Manifolds,'' Bull. Amer. Math. Soc. {\bf 69} (1963) 422-33.

\bibitem{AS}
M. F. Atiyah and I. M. Singer, ``The Index Of Elliptic Operators: I,'' Ann. Math. {\bf 87} (1968) 484-530.

\bibitem{DW}
R. Dijkgraaf and E. Witten, ``Topological Gauge Theories And Group Cohomology,'' Commun. Math. Phys. {\bf 129}
(1990) 393.


\bibitem{ZQH}
Xiao-Liang Qi, T. Hughes, and Shou-Cheng Zhang, ``Topological Field Theory Of Time-Reversal Invariant Insulators,''
arXiv.:0802.3537.


\bibitem{HCR}
C.-T. Hsieh, G. C. Cho, and S. Ryu, ``Global Anomalies On The Surface Of Symmetry-Protected Topological Phases in
$(3+1)$ Dimensions,'' arXiv:1503.01411. 

\bibitem{FCV}
L. Fidkowski, X. Chen, and A. Vishwanath, ``Non-Abelian Topological Order On The
Surface Of A 3d Topological Superconductor From An Exactly Solved Model,''
 Phys. Review {\bf X 3} (2013) 041016. 

\bibitem{WS}
C. Wang and T. Senthil, ``Interacting Fermionic Topological Insulators/Superconductors In Three
Dimensions,'' arXiv:1401.1142,  Phys. Rev. {\bf B89} (2014) 195124.

\bibitem{MFCV}
M. A. Metlitski, L. Fidkowski, X. Chen, and A. Vishwanath, ``Interaction Effects On 3d Topological
Superconductors: Surface Topological Order From Vortex Condensation, the 16-Fold Way, And
Fermionic Kramers Doublets,'' arXiv:1406.3032. 

\bibitem{KitTwo}
A. Kitaev, ``Homotopy-Theoretic Approach to SPT Phases in Action: $\Z_{16}$ Classification of 
Three-Dimensional Superconductors,'' lecture notes available at 
\url{http://www.ipam.ucla.edu/abstract/?tid=12389&pcode=STQ2015}.






\bibitem{Witold}
E. Witten, ``An $\SU(2)$ Anomaly,'' Phys. Lett. {\bf B117} (1982) 324-8.

\bibitem{Witoldtwo}
E. Witten, ``Global Anomalies In String Theory,''  in W. A. Bardeen and A. R. White, eds., 
{\it Symposium On Anomalies, Geometry, Topology} (World Scientific, 1985),
pp. 61-99,
available at \url{https://www.sns.ias.edu/witten/symp-global-anomalies-stringtheory-1985}.


\bibitem{R}
N. Redlich, ``Parity Violation And Gauge Noninvariance of the Effective  Gauge Field Action In Three Dimensions,''  Phys. Rev. {\bf D29} (1984) 2366.

\bibitem{Semenoff}
A. Niemi and G. Semenoff, ``Anomaly-Induced Fermion Fractionization And Effective Actions In Odd-Dimensional Spacetimes,''
Phys. Rev. Lett. {\bf 51} (1983) 2077-80.

\bibitem{ADM}
L. Alvarez-Gaum\'e, S. Della Pietra, and G. Moore, ``Anomalies And Odd Dimensions,'' Ann. Phys. {\bf 163} (1985) 288-317.

\bibitem{Global}
E. Witten, ``Global Gravitational Anomalies,'' Commun. Math. Phys. {\bf 100} (1985) 197-229.

\bibitem{APS}
M. F. Atiyah, I. M. Singer, and V. Patodi, ``Spectral Asymmetry In Riemannian Geometry,'' Math. Proc. Camb. Phil. Soc. {\bf{77} } (1975) 43-69.

\bibitem{DF}
X. Dai and D. S. Freed, ``$\eta$-Invariants And Determinant Lines,'' J. Math. Phys. {\bf 35} (1994) 5155-94.

\bibitem{Dinstanton}
E. Witten, ``World-Sheet Corrections Via D-Instantons,'' JHEP 0002:030 (2000), arXiv:hep-th/9907041.

\bibitem{KTTW}
A. Kapustin, R. Thorngren, A. Turzillo, and Z. Wang, ``Fermionic Symmetry Protected Topological Phases And Cobordisms,''
arXiv:1406.7329.

\bibitem{Freed} 
D. Freed, ``Short Range Entanglement And Invertible Field Theories,'' arXiv.1406:7278.

\bibitem{Bunke}
U. Bunke, ``String Structures And Trivialisations Of A Pfaffian Line Bundle,''  arXiv:0909.0846.

\bibitem{DMW}
E. Diaconescu, G. W. Moore, and E. Witten, ``$E_8$ Gauge Theory, and a Derivation of K-Theory From M-Theory,'' 
Adv. Theor. Math. Phys. {\bf 6} (2003) 1031-1134,
arXiv:hep-th/0005090.

\bibitem{FM}
D. Freed and G. Moore, ``Setting The Quantum Integrand Of M-Theory,'' Commun. Math. Phys. {\bf 263} (2006) 89-132.

\bibitem{MW}
V. Mikhaylov and E. Witten, ``Branes And Supergroups,'' arXiv:1410.1175.

\bibitem{Monnier}
S. Monnier, ``Hamiltonian Anomalies From Extended Field Theories,'' arXiv:1410.7442.

\bibitem{Metlitski}
M. Metlitski, ``$S$-Duality of  $\U(1)$ Gauge Theory With $\theta=\pi$ on Non-Orientable
Manifolds: Applications to Topological Insulators and
Superconductors,'' to appear.

\bibitem{FK}
L. Fidkowski and A. Kitaev, ``The Effects Of Interactions On The Topological Classification Of Free Fermion Systems,''
arXiv:0904.2197.




\bibitem{DCG} C. G. Callan, Jr., R. Dashen, and D. J. Gross, ``Towards A Theory Of The Strong Interactions,'' Phys. Rev. {\bf D17}
(1978) 2717.

\bibitem{Kiskis}
J. Kiskis, ``Fermion Zero-Modes And Level Crossing,'' Phys. Rev. {\bf D18} (1978) 3690.



\bibitem{RF}
C. Rosenberg and M. Franz, ``Witten Effect In A Crystalline Topological Insulator,''
arXiv:1001.3179.

\bibitem{JR}
R. Jackiw and C. Rebbi, ``Solitons  With Fermion Number 1/2,'' Phys. Rev. {\bf D13} (1976) 3398.

\bibitem{Dyons}
E. Witten, ``Dyons Of Charge $e\theta/2\pi$,'' Phys. Lett. {\bf B86} (1979) 283.



\bibitem{AStwo}
M. F. Atiyah and I. M. Singer, ``The Index Of Elliptic Operators: V,''  Ann. Math. {\bf 93} (1971) 139-149.

\bibitem{HW}
P. Horava and E. Witten, ``Heterotic and Type I String Dynamics From Eleven Dimensions,'' Nuel. Phys. {\bf B460} (1996) 506-24, arXiv:hep-th/0510209.

\bibitem{Hsieh}
C.-T. Hsieh, O. M. Sule, G. Y. Cho, S. Tyu, and R. Leigh, ``Symmetry-Protected Topological Phases, Generaliized Laughlin Argument, and Orientifolds,'' 
Phys. Rev. {\bf B90} (2014) 165134.

\bibitem{Berry}
M. Berry, ``Quantal Phase Factors Accompanying Adiabatic Changes,'' Proc. Roy. Soc. {\bf A392} (1984) 45-57.

\bibitem{G}
P. Gilkey, ``The $\eta$-Invariant Of Even-Dimensional $\mathrm{Pin}_c$ Manifolds,'' Adv. Math. 58 (1985) 243-284.

\bibitem{Stolz}
S. Stolz, ``Exotic Structures On 4-Manifolds Detected By
Spectral Invariants,'' Inv. Math. {\bf 94} (1988) 147-162.


\bibitem{BD}
J. D. Bjorken and S. D. Drell, {\it Relativistic Quantum Mechanics} (McGraw Hill, 1964).



\end{thebibliography}
